\section{Estado del arte}
\label{sec:estado-arte}

Durante el proyecto se construy\'o una base de datos bibliogr\'afica estructurada de
\textbf{65 art\'iculos}, cada uno con an\'alisis individual que registra m\'etodo principal y
secundario, tipo de sensor, tipo de fuente, geometr\'ia del arreglo, procesamiento e
inversi\'on empleados, variables estimadas, limitaciones reportadas, conclusiones y
utilidad espec\'ifica para esta tesis. Esta secci\'on sintetiza ese relevamiento.

\subsection{M\'etodos de caracterizaci\'on geot\'ecnica}

Los ensayos se agrupan seg\'un requieran o no acceso al subsuelo:

\begin{table}[h]
\centering
\footnotesize
\caption{Comparaci\'on entre m\'etodos invasivos y no invasivos.}
\label{tab:invasivos}
\begin{tabularx}{\textwidth}{@{}l X X@{}}
\toprule
\textbf{Caracter\'istica} & \textbf{Invasivos (borehole)} & \textbf{No invasivos} \\
\midrule
Ejemplos & Cross-hole (CHT), Down-hole, P-S logging, SPT, CPT, SCPT & SASW, MASW, ReMi, SPAC, HVSR, refracci\'on \\
Costo & Elevado (requiere perforaci\'on) & Considerablemente menor \\
Volumen muestreado & Puntual y localizado & Grandes vol\'umenes representativos \\
Incertidumbre & Baja (medici\'on directa) & Mayor (problema inverso indirecto) \\
Log\'istica & Pesada, equipo de perforaci\'on & Liviana, desplegable por pocas personas \\
Estado del suelo & Perturba el entorno del sondeo & No perturbado \\
\bottomrule
\end{tabularx}
\end{table}

\begin{figure}[ht]
\centering
\includegraphics[width=0.96\textwidth]{comparacion_metodos.png}
\caption{Mapa comparativo de los m\'etodos revisados. Resume qu\'e magnitud aporta cada
familia, la log\'istica que exige y el lugar que ocupa MASW activo dentro del proyecto;
no representa una jerarqu\'ia universal de exactitud.}
\label{fig:comparacion-metodos}
\end{figure}

La referencia cuantitativa m\'as s\'olida sobre la confiabilidad relativa de ambas familias
es el proyecto \textbf{InterPACIFIC}, un ejercicio de comparaci\'on entre laboratorios con
14 equipos participantes. Su primera parte~\cite{Garofalo2016a} cuantific\'o la dispersi\'on
\emph{entre} distintos operadores aplicando m\'etodos de ondas superficiales al mismo
sitio; la segunda~\cite{Garofalo2016b} compar\'o esos resultados contra m\'etodos de
perforaci\'on. Las gu\'ias de buenas pr\'acticas derivadas del proyecto~\cite{Foti2018}
constituyen la referencia metodol\'ogica que orienta el dise\~no experimental de este
trabajo; no fijan una norma \'unica aplicable a toda configuraci\'on.
Dentro del intervalo de longitudes de onda considerado confiable, el ejercicio encontr\'o
coeficientes de variaci\'on entre analistas de aproximadamente \SIrange{5}{10}{\percent}
para la velocidad de fase media, mientras que las incertidumbres declaradas por los
propios analistas estuvieron t\'ipicamente entre \SIrange{5}{20}{\percent}. La segunda
parte mostr\'o que los m\'etodos invasivos resuelven mejor capas delgadas y contrastes
abruptos, pero que ambas familias producen valores de $V_{S,30}$ comparables.

En la misma l\'inea, Stephenson~\cite{Stephenson2005} compar\'o a ciegas ReMi y MASW con
cuatro sondeos registrados hasta al menos \SI{260}{\meter}. Los promedios de $V_S$ hasta
30, 50 y \SI{100}{\meter} quedaron generalmente dentro de \SI{15}{\percent} de la
referencia; en dos sondeos el promedio a \SI{100}{\meter} difiri\'o menos de
\SI{3}{\percent}, y la amplificaci\'on espectral entre 1 y \SI{10}{\hertz} qued\'o dentro
de \SI{15}{\percent}. Ninguno de los dos m\'etodos fue consistentemente superior.
Comina \emph{et al.}~\cite{Comina2011}, sobre ocho sitios, reportaron repetibilidad de
medici\'on t\'ipicamente mejor que \SI{5}{\percent} y variaci\'on de $V_{S,30}$ cercana a
\SI{2}{\percent} entre modelos equivalentes cuando la banda experimental cubre las
longitudes de onda necesarias. Con una longitud de onda m\'axima superior a
\SI{60}{\meter}, los resultados se aproximaron a los invasivos; el estimador simplificado
basado en $c_R$ a \SI{36}{\meter} tuvo una dispersi\'on cercana a \SI{6}{\percent} y una
desviaci\'on RMS de aproximadamente \SI{7}{\percent}, pero no sustituye la inversi\'on.

\subsubsection{Qu\'e se lleva este trabajo de InterPACIFIC}
\label{sec:buenas-practicas}

El hallazgo estructural de InterPACIFIC, m\'as importante que cualquiera de sus n\'umeros
individuales, es la \textbf{separaci\'on entre dos fuentes de error que suelen confundirse}:

\begin{itemize}[itemsep=0.2em]
    \item La \textbf{curva de dispersi\'on experimental} presenta variabilidad
    \emph{baja} entre equipos independientes que miden el mismo sitio. Es decir, la etapa
    de \emph{medici\'on} ---instrumento, arreglo, procesamiento hasta la imagen de
    dispersi\'on--- est\'a esencialmente resuelta y es reproducible.
    \item La dispersi\'on entre los \textbf{perfiles $V_S$} finales es sustancialmente
    mayor, y su origen dominante es la \textbf{parametrizaci\'on del modelo y la no
    unicidad de la inversi\'on} (\S\ref{sec:no-unicidad}), no el ruido de medici\'on.
\end{itemize}

La lectura para este proyecto es directa y conviene enunciarla sin adornos: \textbf{el
cuello de botella de la exactitud no est\'a en el instrumento que este trabajo construye,
sino en la inversi\'on}. Eso no invalida el aporte ---un instrumento que no mide no habilita
ninguna inversi\'on--- pero s\'i fija expectativas: mejorar el \emph{front-end} mejora la
banda \'util y la SNR de la curva de dispersi\'on, que es exactamente la etapa que
InterPACIFIC encuentra reproducible; no reduce por s\'i mismo la incertidumbre del perfil
final.

De las gu\'ias de buenas pr\'acticas~\cite{Foti2018} derivadas del proyecto, este trabajo
adopta expl\'icitamente las siguientes, y en cada caso se indica d\'onde se materializan:

\begin{center}\footnotesize
\begin{tabularx}{\textwidth}{@{}l X l@{}}
\toprule
Buena pr\'actica & C\'omo se implementa aqu\'i & D\'onde \\
\midrule
Apilamiento de m\'ultiples golpes & Promediado por punto de fuente, con rechazo expl\'icito de capturas an\'omalas y compuerta \texttt{reviewed} por traza & \S\ref{sec:metodologia} \\
Control del \emph{offset} m\'inimo fuente--primer receptor & Registro del \emph{offset} por nodo en la interfaz; criterio de campo contra efecto de campo cercano & \S\ref{sec:estado-arte-arreglo} \\
Declaraci\'on de la banda v\'alida de la curva & Selecci\'on manual de la regi\'on de dispersi\'on con restricci\'on anti-\emph{aliasing} $c\ge 2\,\Delta x\,f$ en el \emph{picking} & \S\ref{sec:arq-interfaces} \\
Reporte de incertidumbre, no de un \'unico perfil & Inversi\'on evolutiva que entrega poblaci\'on de modelos, no una \'unica soluci\'on & \S\ref{sec:resultados} \\
Documentaci\'on de la geometr\'ia real, no la planificada & Geometr\'ia registrada por captura en los metadatos & Ap\'endice~\ref{sec:inventario} \\
Verificaci\'on de la respuesta instrumental & Caracterizaci\'on de la cadena por barrido y contraste con el modelo & \S\ref{sec:metodologia} \\
\bottomrule
\end{tabularx}
\end{center}

Del trabajo de Comina \emph{et al.}~\cite{Comina2011} se retiene adem\'as un
\textbf{criterio de dise\~no directamente accionable}: si la banda experimental alcanza
longitudes de onda m\'aximas superiores a \SI{60}{\meter}, la no unicidad deja de afectar
significativamente la exactitud de $V_{S,30}$, con incertidumbres del orden del
\SIrange{5}{6}{\percent}. Ese umbral de $\lambda_{\max}>\SI{60}{\meter}$ es
\textbf{exactamente la especificaci\'on que la compensaci\'on del transductor
persigue habilitar} (\S\ref{sec:geo-park}) y el criterio contra el cual debe contrastarse
la apertura del arreglo desplegado (\S\ref{sec:estado-arte-arreglo}).

\subsection{Evoluci\'on de los m\'etodos de ondas superficiales}

\subsubsection{M\'etodos activos}

\begin{itemize}[itemsep=0.25em]
    \item \textbf{Steady-State Rayleigh Method} (d\'ecada de 1950): fuente arm\'onica de
    frecuencia variable y b\'usqueda manual de puntos en fase. Conceptualmente directo pero
    extremadamente lento.

    \item \textbf{SASW --- Spectral Analysis of Surface Waves}~\cite{Nazarian1984}: dos
    receptores y an\'alisis del espectro cruzado para obtener el desfasaje en funci\'on de la
    frecuencia. Introduce la funci\'on de coherencia como criterio de calidad de la
    medici\'on. Su debilidad estructural es el \emph{phase unwrapping}: la fase medida est\'a
    envuelta en $(-180^{\circ},+180^{\circ}]$ y los errores en la identificaci\'on de los
    saltos se propagan directamente a la curva de dispersi\'on. Requiere adem\'as m\'ultiples
    repeticiones con distintos espaciamientos, lo que penaliza la productividad de campo.

    \item \textbf{MASW --- Multichannel Analysis of Surface Waves}~\cite{Park1999}: arreglo
    multicanal registrado simult\'aneamente y transformado a un dominio donde la dispersi\'on
    es observable. Al emplear todos los canales a la vez, la transformaci\'on act\'ua como un
    filtro espacial que separa modos y discrimina ruido incoherente, eliminando el
    problema de \emph{unwrapping} y mejorando dr\'asticamente la tasa de producci\'on de
    campo. Es el m\'etodo adoptado en este proyecto.
\end{itemize}

La tr\'iada fundacional del MASW la completan el trabajo de
procesamiento~\cite{Park1998}, que introduce el m\'etodo de \emph{phase-shift} ---el
est\'andar de facto implementado en pr\'acticamente todos los programas de an\'alisis MASW,
libres y comerciales---, y el de inversi\'on~\cite{Xia1999}, que define el algoritmo
iterativo basado en el Jacobiano de la velocidad de fase respecto de $V_S$ por capa.
Miller \emph{et al.}~\cite{Miller1999} demostraron tempranamente la aplicaci\'on al mapeo
de basamento rocoso, y el trabajo de revisi\'on de Socco y Strobbia~\cite{Socco2004}
sistematiz\'o la metodolog\'ia completa.

\subsubsection{M\'etodos pasivos (evaluados, no adoptados)}

Los m\'etodos pasivos aprovechan el ruido s\'ismico ambiental en lugar de una fuente
activa, lo que elimina el problema de generar energ\'ia s\'ismica repetible:

\begin{itemize}[itemsep=0.2em]
    \item \textbf{SPAC}~\cite{Aki1957}: el trabajo fundacional de los m\'etodos pasivos,
    basado en la autocorrelaci\'on espacial de ondas estacionarias estoc\'asticas registradas
    en arreglos circulares.
    \item \textbf{ReMi}~\cite{Louie2001}: emplea el mismo arreglo lineal del MASW activo,
    aprovechando el ruido de tr\'ansito. Su factibilidad econ\'omica es muy alta ---no
    requiere fuente--- pero su resoluci\'on depende de la direccionalidad del campo de
    ruido.
    \item \textbf{HVSR}~\cite{Nakamura1989,Lermo1993}: relaci\'on espectral
    horizontal/vertical con una sola estaci\'on tri-axial. Es el m\'etodo de mejor relaci\'on
    costo--beneficio para estimar la frecuencia fundamental del sitio, aunque no entrega
    directamente el perfil $V_S$. Las gu\'ias SESAME~\cite{SESAME2004} sistematizan su
    aplicaci\'on.
\end{itemize}

Park \emph{et al.}~\cite{Park2007} muestran que activo y pasivo pueden combinarse con el
mismo equipamiento, sin costo adicional de hardware: el activo cubre las frecuencias
altas y el pasivo extiende la curva hacia frecuencias bajas, es decir, hacia mayores
profundidades.

\nota{\textbf{Oportunidad}: el sistema desarrollado puede, en principio, adquirir en modo
pasivo sin ninguna modificaci\'on de hardware ---basta capturar sin disparar la fuente,
y la capacidad de captura continua de 10 minutos validada ya lo permite. Es una
extensi\'on de bajo costo que ampliar\'ia el alcance del trabajo y que conviene mencionar
en trabajos futuros. ¿Se consider\'o?}

\subsection{Dise\~no del arreglo: los criterios cuantitativos}
\label{sec:estado-arte-arreglo}

Esta es la parte del estado del arte que impacta m\'as directamente sobre las decisiones
de campo, y sobre la que existe literatura espec\'ifica:

\begin{itemize}[itemsep=0.25em]
    \item \textbf{Par\'ametros \'optimos de campo}: Park \emph{et al.}~\cite{Park2002}
    sistematizan la elecci\'on de espaciamiento, longitud de arreglo y \emph{offset} en
    funci\'on de la profundidad objetivo.
    \item \textbf{Efecto de campo cercano}: a distancias menores a aproximadamente media
    longitud de onda, el campo no est\'a dominado por la onda superficial plana y la
    velocidad de fase medida queda sesgada. Yoon y Rix~\cite{Yoon2009} cuantifican el
    efecto para m\'etodos de arreglo con fuente activa.
    \item \textbf{Offset m\'inimo}: Xu \emph{et al.}~\cite{Xu2006} entregan una estimaci\'on
    cuantitativa de la distancia m\'inima fuente--primer receptor, complemento directo del
    anterior.
    \item \textbf{Ecuaciones de dise\~no}: Xia \emph{et al.}~\cite{Xia2006} proveen
    relaciones simples que vinculan los par\'ametros del arreglo con el rango de
    frecuencias y la profundidad de investigaci\'on alcanzables.
    \item \textbf{Apilamiento}: Neducza~\cite{Neducza2007} analiza la suma coherente de
    m\'ultiples disparos, que mejora la relaci\'on se\~nal--ruido en factor $\sqrt{N}$ y es
    la pr\'actica adoptada en este proyecto (m\'ultiples repeticiones por punto de fuente).
    \item \textbf{An\'alisis CMP}: Hayashi y Suzuki~\cite{Hayashi2004} introducen la
    correlaci\'on cruzada por punto medio com\'un, que permite construir secciones
    pseudo-2D de $V_S$ a partir de adquisici\'on rodante.
\end{itemize}

Estos criterios (\cite{Park2002,Yoon2009,Xu2006,Xia2006}) se aplican en la
Secci\'on~\ref{sec:resultados} a la geometr\'ia realmente usada. La conclusi\'on no se
deduce s\'olo de la frecuencia m\'inima observada: exige considerar simult\'aneamente
espaciamiento, apertura, \emph{offset}, energ\'ia de fuente y banda con SNR suficiente.

\subsection{Modos superiores e inversi\'on}

\subsubsection{El problema de los modos}

Tokimatsu \emph{et al.}~\cite{Tokimatsu1992} caracterizaron los efectos de los modos
superiores sobre las caracter\'isticas de dispersi\'on, mostrando que en perfiles con
inversiones de velocidad la energ\'ia puede migrar a modos superiores y la curva medida
deja de corresponder al modo fundamental. Buchen y Ben-Hador~\cite{Buchen1996}
sistematizaron el c\'alculo de modos libres, base de los motores de modelado directo
actuales.

El tratamiento moderno explota los modos superiores en lugar de evitarlos:
Xia \emph{et al.}~\cite{Xia2003} invierten conjuntamente modo fundamental y superiores;
Luo \emph{et al.}~\cite{Luo2009} separan modos mediante transformada de Radon lineal de
alta resoluci\'on; y Maraschini y Foti~\cite{Maraschini2010} proponen una inversi\'on
multimodal por Monte Carlo, aplicada a sitios con basamento somero por Bergamo
\emph{et al.}~\cite{Bergamo2011}.

El \emph{pipeline} desarrollado permite separar varias curvas mediante regiones poligonales
sobre la imagen de dispersi\'on y dispone de motores capaces de invertirlas en conjunto.
Esa es una \textbf{capacidad de la infraestructura}, no la descripci\'on del resultado de
campo que se conserva actualmente: la inversi\'on persistida usa s\'olo el modo~0 asignado
por el operador, sin identificaci\'on modal independiente,
con el \emph{port} num\'erico de MASWavesPy. La comparaci\'on multimodal entre motores queda
como validaci\'on futura.

\paragraph{Qu\'e se eligi\'o, concretamente.} Conviene ser expl\'icito porque el p\'arrafo
anterior describe una capacidad y no una elecci\'on. La configuraci\'on adoptada es:

\begin{itemize}[itemsep=0.15em]
    \item \textbf{Extracci\'on de curvas}: selecci\'on manual de \emph{regiones
    poligonales} sobre la imagen de dispersi\'on, una regi\'on por modo, produciendo $N$
    curvas simult\'aneas. La alternativa autom\'atica ---seguimiento de m\'aximos--- se
    descart\'o porque en presencia de salto modal el m\'aximo cambia de modo sin aviso, que
    es exactamente el modo de falla que Tokimatsu \emph{et al.}~\cite{Tokimatsu1992}
    describen.
    \item \textbf{Modelado directo}: \textbf{disba}, implementaci\'on por matriz de
    propagaci\'on compilada con Numba.
    \item \textbf{Inversi\'on}: \textbf{evodcinv}, b\'usqueda evolutiva sobre disba, capaz
    de invertir varios modos en conjunto y de entregar una poblaci\'on de modelos en vez de
    una \'unica soluci\'on ---lo que se alinea con la buena pr\'actica de reportar
    incertidumbre (\S\ref{sec:buenas-practicas}).
    \item \textbf{Motor del estado persistido de campo} que se reporta en resultados:
    \emph{port} num\'erico de MASWavesPy sobre el modo~0.
\end{itemize}

\paragraph{Por qu\'e esa elecci\'on, y no otra.} La justificaci\'on es deliberadamente
simple y conviene declararla como tal en lugar de construir un argumento t\'ecnico
\emph{a posteriori}: \textbf{se apost\'o por implementaciones ya validadas por la
comunidad \emph{open-source}} en lugar de reimplementar el modelado directo. Un motor de
matriz de propagaci\'on es c\'odigo delicado ---inestabilidad num\'erica a alta frecuencia,
p\'erdida de precisi\'on en capas delgadas, b\'usqueda de ra\'ices que puede saltarse
modos--- y reescribirlo habr\'ia consumido meses sin aportar nada nuevo, adem\'as de
introducir una fuente de error imposible de separar de los errores del instrumento que
s\'i se quiere evaluar. \textbf{El aporte de este proyecto est\'a en la instrumentaci\'on y
en la cadena de adquisici\'on, no en el motor de inversi\'on}; usar herramientas
consolidadas para la parte que no es el aporte es la decisi\'on econ\'omicamente correcta
y adem\'as la que hace los resultados reproducibles por terceros. Los criterios
operativos que decantaron la elecci\'on concreta fueron tres: integrabilidad directa en
Python (el resto del \emph{pipeline} ya lo es), licencia permisiva, y soporte nativo de
inversi\'on multimodal.

\subsubsection{No unicidad y regularizaci\'on}
\label{sec:no-unicidad}

La inversi\'on es un problema mal condicionado: distintos perfiles $V_S(z)$ pueden producir
curvas de dispersi\'on indistinguibles dentro del error experimental. Foti \emph{et
al.}~\cite{Foti2009} analizan expl\'icitamente las consecuencias de esa no unicidad sobre
el an\'alisis de respuesta de sitio ---es decir, sobre el uso final del perfil.

Las estrategias de mitigaci\'on relevadas incluyen:

\begin{itemize}[itemsep=0.15em]
    \item \textbf{Parametrizaci\'on sistem\'atica}: Cox y Teague~\cite{Cox2016} proponen el
    m\'etodo de \emph{layering ratios} para elegir el n\'umero y espesor de capas en ausencia
    de informaci\'on a priori ---un problema pr\'actico recurrente.
    \item \textbf{B\'usqueda global}: Wathelet \emph{et al.}~\cite{Wathelet2004} aplican un
    algoritmo de b\'usqueda directa (algoritmo de vecindad, base de Dinver/Geopsy) que
    explora el espacio de modelos sin depender del modelo inicial, a diferencia de los
    m\'etodos locales tipo Gauss-Newton.
    \item \textbf{Cuantificaci\'on de incertidumbre}: Griffiths \emph{et
    al.}~\cite{Griffiths2016} mapean el desajuste de la curva de dispersi\'on y la
    incertidumbre del perfil hacia la variabilidad de la respuesta de sitio.
\end{itemize}

La inversi\'on requiere adem\'as fijar par\'ametros auxiliares ---densidad y raz\'on de
Poisson por capa--- de influencia limitada pero no nula. En particular, ignorar el nivel
fre\'atico equivale a asumir una raz\'on de Poisson incorrecta para las capas saturadas
(donde el comportamiento no drenado lleva $\nu\to0{,}5$), con consecuencias
significativas sobre el perfil resultante.

\subsection{Instrumentaci\'on: ge\'ofonos y adquisici\'on}

El ge\'ofono de bobina m\'ovil sigue siendo el transductor dominante por robustez, bajo
costo y ausencia de necesidad de alimentaci\'on. Su limitaci\'on ---la ca\'ida de respuesta
por debajo de la frecuencia natural, Secci\'on~\ref{sec:geofono}--- motiva una l\'inea de
trabajo activa: Ma \emph{et al.}~\cite{Ma2023} proponen el m\'etodo de extensi\'on por
aumento del amortiguamiento que es la base de la cadena implementada en este proyecto
(desarrollado en detalle en la Secci\'on~\ref{sec:ma2023}), y Deng \emph{et
al.}~\cite{Deng2025} reportan un circuito de acondicionamiento de baja frecuencia que
reduce la frecuencia de corte de un ge\'ofono de \SI{4.5}{\hertz} hasta
\SI{0.16}{\hertz}.

En cuanto a plataformas de adquisici\'on de bajo costo existen antecedentes de interfaces
embebidas para ge\'ofonos~\cite{Kafadar2020,Wijayaraja2024} y de nodos reconfigurables
para redes de sensores~\cite{Guevara2015}. La tendencia general es reemplazar el
sismógrafo centralizado con cableado multicanal por nodos aut\'onomos inal\'ambricos, lo
que reduce dr\'asticamente la log\'istica de despliegue a costa de introducir el problema
de la sincronizaci\'on temporal distribuida.

\subsubsection{Plataformas mixtas programables y asistencia digital}

Las plataformas de se\~nal mixta programables como el PSoC~5LP~\cite{InfineonPSoC}
integran bloques anal\'ogicos configurables, conversores, filtros digitales programables,
DMA y procesador en un \'unico dispositivo. Sus bloques anal\'ogicos ofrecen, sin embargo,
menor precisi\'on, mayores \emph{offsets} y peor apareamiento que circuitos dedicados.

Frente a ello, el enfoque de \emph{digitally assisted analog design} de Murmann y
Boser~\cite{Murmann2004,Murmann2006} propone usar recursos digitales para asistir al
circuito anal\'ogico en lugar de mejorar los componentes. Aunque se investig\'o
principalmente a nivel de circuito integrado, este proyecto extiende la filosof\'ia al
nivel de instrumentaci\'on de bajo costo.

Las t\'ecnicas cl\'asicas de compensaci\'on de \emph{offset} siguen tres caminos:
correcci\'on num\'erica posterior a la adquisici\'on~\cite{InfineonAN68403,Lopin2018};
ajuste f\'isico de cada etapa~\cite{Mohan2020}; y \emph{autozero} o estabilizaci\'on por
\emph{chopper}~\cite{EnzTemes1996,Liu2020}. Los servos de continua continuos introducen
adem\'as un compromiso inherente entre tiempo de establecimiento y ancho de
banda~\cite{Liu2020,Antoniadis2025}, y los \emph{front-ends} de precisi\'on que operan
cerca del l\'imite de ruido t\'ermico del sensor~\cite{Kirchhoff2017} representan el
enfoque opuesto: mejorar el componente anal\'ogico en lugar de asistirlo. El aporte de
este proyecto se ubica entre ambos.

\subsection{Herramientas de procesamiento e inversi\'on}

\begin{itemize}[itemsep=0.15em]
    \item \textbf{MASWavesPy}~\cite{OlafsdottirNGM2024,Olafsdottir2018,Olafsdottir2020}:
    an\'alisis MASW e inversi\'on con cuantificaci\'on de incertidumbre por Monte Carlo.
    Integrado como subm\'odulo del proyecto.
    \item \textbf{Geopsy}~\cite{Wathelet2020}: referencia de facto en m\'etodos pasivos y
    HVSR; implementa el algoritmo de b\'usqueda de~\cite{Wathelet2004}.
    \item \textbf{SWprocess}~\cite{Vantassel2022}: orientado a estimaciones robustas de
    la incertidumbre de la curva de dispersi\'on.
    \item \textbf{disba} y \textbf{evodcinv}: modelado directo por matriz de propagaci\'on
    e inversi\'on por algoritmos evolutivos, disponibles como \emph{backend} multimodal del
    proyecto; no son el motor del estado de inversi\'on persistido que se reporta en
    resultados.
\end{itemize}

\paragraph{La comparaci\'on entre implementaciones queda como trabajo futuro.} Debe
quedar constancia expl\'icita de un l\'imite del estado actual: \textbf{la elecci\'on
entre estas herramientas fue pragm\'atica, no el resultado de un \emph{benchmark}}. Se
decidi\'o por integrabilidad en Python, licencia y soporte multimodal
(\S\ref{sec:no-unicidad}), no porque se haya demostrado que alguna produce mejores
perfiles sobre los datos de este proyecto.

La evaluaci\'on sistem\'atica de las distintas implementaciones y m\'etodos propuestos
---MASWavesPy frente a Geopsy frente a SWprocess frente a disba+evodcinv, y dentro de cada
uno las variantes de transformaci\'on al dominio de dispersi\'on (desplazamiento de fase,
$f$--$k$, $\tau$--$p$/Radon) y de estrategia de inversi\'on (b\'usqueda local, algoritmo de
vecindad, evolutiva, Monte Carlo)--- \textbf{queda como trabajo futuro}
(\S\ref{sec:futuros}). Es una comparaci\'on que este proyecto est\'a en condiciones
inusualmente buenas de hacer, porque dispone de algo que la mayor\'ia de los trabajos de
comparaci\'on no tiene: \emph{datos propios de un sitio medido repetidamente con un
instrumento cuya respuesta se caracteriz\'o de forma independiente}. Corriendo todos los
motores sobre exactamente los mismos registros, la dispersi\'on entre perfiles resultantes
ser\'ia atribuible \'unicamente al m\'etodo ---lo que constituye, a escala reducida, el
mismo experimento que InterPACIFIC realiz\'o entre equipos
(\S\ref{sec:buenas-practicas}).

\subsection{Aplicaci\'on final: del perfil \texorpdfstring{$V_S$}{VS} al par\'ametro geot\'ecnico}

El perfil $V_S$ no es un fin en s\'i mismo. Sus usos establecidos son:

\begin{itemize}[itemsep=0.15em]
    \item \textbf{Clasificaci\'on s\'ismica de sitio} mediante $V_{S,30}$, par\'ametro
    normativo en los c\'odigos de dise\~no; Boore~\cite{Boore2004} sistematiza su estimaci\'on
    a partir de modelos de velocidad someros.
    \item \textbf{Evaluaci\'on del potencial de licuefacci\'on}: Andrus y
    Stokoe~\cite{Andrus2000} establecieron la correlaci\'on entre resistencia a la
    licuefacci\'on y velocidad de onda de corte.
    \item \textbf{Correlaci\'on con ensayos de penetraci\'on}, que es la validaci\'on prevista
    para este proyecto y se trata a continuaci\'on.
\end{itemize}

\subsubsection{Correlaciones \texorpdfstring{$V_S$}{VS}--SPT}

Existe un cuerpo consolidado de trabajos que correlacionan $V_S$ obtenido por ondas
superficiales con la resistencia a la penetraci\'on $N$ del SPT. Uma Maheswari \emph{et
al.}~\cite{UmaMaheswari2010} desarrollaron correlaciones para los suelos de Chennai a
partir de 30 sitios urbanos y 200 pares $V_S$--$N$, validadas contra cross-hole;
Dikmen~\cite{Dikmen2009} present\'o correlaciones estad\'isticas por tipo de suelo; y
Alhuay-Le\'on y Trejo-Nore\~na~\cite{Alhuay2021} las establecieron para las arenas
e\'olicas de Lima ---el antecedente regional m\'as cercano metodol\'ogicamente a la
validaci\'on que este proyecto planea realizar.

\subsection{Antecedentes regionales}

Se relevaron espec\'ificamente trabajos de aplicaci\'on en contextos comparables al local,
tanto por tipo de suelo como por disponibilidad de recursos:

\begin{center}\footnotesize
\begin{tabularx}{\textwidth}{@{}l l X@{}}
\toprule
Trabajo & Contexto & Relevancia para este proyecto \\
\midrule
Moffat \emph{et al.}~\cite{Moffat2016} & Chile & Compara $V_{S,30}$ por downhole, MASW y ReMi; 12 ge\'ofonos + maza + Geopsy: configuraci\'on econ\'omica y replicable \\
Alhuay-Le\'on \emph{et al.}~\cite{Alhuay2021} & Per\'u & MASW + SPT sobre arenas e\'olicas; el contraste que este proyecto planea replicar \\
Moya-Guti\'errez \emph{et al.}~\cite{Moya2020} & Colombia & Prospecci\'on geof\'isica y geot\'ecnica combinada; ge\'ofonos de \SI{14}{\hertz} \\
Lima J\'unior \emph{et al.}~\cite{LimaJunior2012} & Brasil & Clima tropical h\'umedo ($>\SI{2000}{\milli\meter}$/a\~no), condici\'on similar a la local; aplicaci\'on a laderas inestables \\
Ayele \emph{et al.}~\cite{Ayele2022} & Etiop\'ia & MASW con equipamiento est\'andar en contexto de recursos limitados; demuestra viabilidad \\
\bottomrule
\end{tabularx}
\end{center}

Los sitios de investigaci\'on geot\'ecnica europeos~\cite{Long2007,Blaker2019,Quinteros2019}
se relevaron como marco de comparaci\'on metodol\'ogica, por disponer de caracterizaci\'on
exhaustiva de referencia.

\subsection{Brecha identificada y posicionamiento del trabajo}

Del relevamiento se desprende que la teor\'ia MASW, los criterios de dise\~no de arreglo,
las buenas pr\'acticas de adquisici\'on y las herramientas de procesamiento est\'an
disponibles, son maduras y en su mayor\'ia abiertas. La barrera de entrada real para su
aplicaci\'on local es el \textbf{instrumental de adquisici\'on multicanal sincronizado}.

Los sistemas comerciales son costosos y cerrados. Los desarrollos acad\'emicos de bajo
costo publicados~\cite{Kafadar2020,Wijayaraja2024} suelen limitarse a la interfaz de un
canal o a la demostraci\'on del \emph{front-end}, sin abordar el sistema completo
---sincronizaci\'on inal\'ambrica multi-nodo, almacenamiento, interfaz de operador en campo
y \emph{pipeline} de procesamiento--- como un conjunto integrado y validado. Y en el eje
de la instrumentaci\'on anal\'ogica, los trabajos de extensi\'on de respuesta de baja
frecuencia~\cite{Ma2023,Deng2025} se validan sobre un canal de laboratorio, no sobre un
arreglo desplegado en campo.

Este proyecto se posiciona en esa brecha: integrar un prototipo de extremo a extremo con
componentes de bajo costo y validar progresivamente sus subsistemas. El aporte
metodol\'ogico espec\'ifico ---el \emph{front-end} anal\'ogico asistido digitalmente con
auto-calibraci\'on--- surgi\'o como respuesta a un problema concreto encontrado durante la
construcci\'on, y fue sometido al congreso URUCON~2026; el env\'io no se contabiliza como
revisi\'on externa recibida hasta contar con un dictamen.

\nota{Esta secci\'on es el n\'ucleo formal de la Primera Presentaci\'on (el Art.~19 pide
espec\'ificamente un ``documento de Estado del Arte''). En la versi\'on de 15 p\'aginas
deber\'ia \emph{crecer} en peso relativo respecto de las secciones de implementaci\'on:
sugerido 4--5 p\'aginas, conservando \S3.1, \S3.2, \S3.3, \S3.5.1, \S3.8 y \S3.9 completos,
y comprimiendo \S3.4 y \S3.6.}
